% OpenAI mathematical source adc7f1241b42e322a6451854ab7e4b4c146bf78a
% Lines 130–147 followed by 255–281.
\eqref{comparator:typical-set}, and a mean tree with metrics
\eqref{comparator:metrics}. There is a sequence of common auxiliary
labels $\lambda_k$ and nonzero vectors
\begin{equation}
 \mathrm{pre}_k=\pi_R^{\lambda_k}
  (\widetilde\Omega\otimes\mathrm{Bell}_{CR})^{\otimes k}
       \in\mathcal S_k,
 \qquad \|\mathrm{pre}_k\|\le1,
 \label{comparator:prevector}
\end{equation}
Here $\pi_R^{\lambda_k}$ is the Schur-label projection on the $R$ copies.
These vectors are physical mean-energy eigenvectors of eigenvalue
$\widetilde E_0$ and have label $\lambda_k$ on both $C$ and $R$.
For every tree above, set
\[
 N=\|\mathcal M^{-s}\mathrm{pre}_k\|,
 \qquad v=\mathcal M^{-s}\mathrm{pre}_k/N,
 \qquad
The auxiliary label must have at least inverse polynomial probability in
$\mathrm{post}_k$.  We therefore select it using $\psi$, then impose
it on the initial auxiliary pair.  Its initial probability need only
be nonzero; the prevector remains unnormalized.

In $\psi^{\otimes k}$, the iid surprisal on $R$ is
$kS'+o(k)$ in probability.  For instance, at fixed system Chebyshev's
inequality gives a window of width $k^{3/4}$.  By
\eqref{comparator:label-remainder}, the difference between that surprisal
and $F_R$ exceeds $k^{3/4}$ with probability at most
$\poly(k)e^{-k^{3/4}/2}$.  The number of labels is polynomial.
Consequently one can select $\lambda_k$ such that
\begin{equation}
 q_k:=\|\pi_R^{\lambda_k}\psi^{\otimes k}\|^2
       \ge\poly(k)^{-1},\qquad
 \log d_{\lambda_k}=kS'+o(k).
 \label{comparator:high-label}
\end{equation}
The same label has positive probability on the initial uniform pair:
every label allowed for $d_*$ occurs in the tensor power of a full-rank
maximally mixed density.  This proves that \eqref{comparator:prevector}
is nonzero.  Its physical part is unchanged by the auxiliary projection,
so it is an exact mean-energy eigenvector.  On the repeated uniform pair,
the $C$ and $R$ labels agree.  Each metric preserves these two labels:
a central permutation operator on an auxiliary system commutes with the
simultaneous permutations on a larger region containing it.  The tree
preserves them as well.
